\chapter{Background}

[\textit{Must be present in FYDP-1 Report and also in Final Report}]

Every chapter should start with 1-2 sentences on the outline of the chapter.

\section{Preliminaries}
In this section, you have to provide the necessary background knowledge to understand the rest of the report \cite{kleinberg2006algorithm}.
\section{Literature Review}
The Literature Review section summarizes and analyzes existing research, systems, and technologies related to the project, showing what has already been done and where your work fits in. It typically includes a review of relevant papers, methods, tools, and findings, along with a brief comparison of their strengths and limitations. The goal is not just to describe past work, but to critically analyze it and identify trends, gaps, and opportunities that justify your proposed solution.

In a thesis (research) work, for example, a student working on an AI-based disease prediction model may review several studies using CNN, SVM, or hybrid models, compare their accuracy (e.g., 80–90\%), and discuss limitations such as small datasets or lack of generalization. In a project (system/software) work, such as a web-based management system, the review may include existing applications or frameworks, discussing features, usability, and shortcomings like lack of scalability or poor user interface. Overall, this section demonstrates understanding of the domain, awareness of prior work, and the need for the proposed solution.

This section will contain you literature review.
\subsection{Similar Applications}

The Similar Applications section describes existing real-world systems, software, or applications that are closely related to your project. It focuses on practical implementations (not just research papers) and explains their features, technologies, strengths, and limitations. The purpose is to show awareness of current solutions and to highlight how your system is different or improved.

In a thesis (research) work, for example, if the topic is AI-based disease prediction, this section may discuss existing applications like hospital decision-support systems or diagnostic apps, describing their capabilities and limitations (e.g., limited accuracy or lack of local data). In a project (system/software) work, such as a student management system, similar applications may include existing university portals or commercial software, analyzing features like record management, reporting, and user interface, while noting gaps such as poor usability or lack of customization. Overall, this section demonstrates practical relevance, comparison with real systems, and justification for the proposed solution.

\subsection{Related Research}
The Related Research section reviews and discusses previous academic studies, journal papers, and conference works that are directly connected to your thesis or project topic. It focuses on research methods, models, algorithms, datasets, and key findings, showing how other researchers have approached similar problems. The section should not only summarize these works but also compare them, highlight their strengths and limitations, and identify research gaps that your work aims to address.

In a thesis (research) work, for example, a student working on an AI-based disease prediction system may analyze several studies using techniques like CNN, SVM, or hybrid models, comparing their performance (e.g., accuracy, PR-AUC) and noting limitations such as small datasets or lack of generalization. In a project (system/software) work, such as a smart traffic system, the section may include research on traffic optimization algorithms or IoT-based solutions, discussing their effectiveness and constraints. Overall, this section demonstrates deep understanding of the research landscape and positions your work within existing scientific contributions.

\section{Gap Analysis}
Gap analysis identifies the limitations, shortcomings, or unexplored areas in existing research or systems and explains how the proposed work aims to address those gaps.
Recommended Table Format. The table columns will include: Existing Work, Limitations, Identified Gap, and Proposed Solution.
Here summarise the gap where you intend to work.
\section{Summary}
